\documentclass[12pt, a4paper]{article}

% ===== ESSENTIAL PACKAGES =====
\usepackage{geometry}
\geometry{a4paper, margin=1in}
\usepackage{setspace}
\onehalfspacing
\usepackage{hyperref}
\usepackage{amsthm}
\newtheorem{definition}{Definisjon}
\usepackage{catchfile}

% ===== CUSTOM ENVIRONMENTS =====
\newenvironment{argument}{\begin{quote}\itshape\textbf{Argument: }}{\end{quote}}

% ===== DOCUMENT BEGIN =====
\title{Exphil03 Obligatorisk oppgave nr. 2}
\author{Oliver Ekeberg}
\date{\today}

\begin{document}
\maketitle

\tableofcontents

\section{Innledning}

René Descartes (1596–1650) regnes som grunnleggeren av substansdualismen – læren om at mennesket består av to vesensforskjellige substanser: en tenkende sjel (\textit{res cogitans}) og en utstrakt kropp (\textit{res extensa}). Spørsmålet om hvordan disse to kan virke sammen, står sentralt i brevene mellom Descartes og prinsesse Elisabeth av Böhmen (1596–1662). 

I brevet 6. mai 1643 stiller Elisabeth et grunnleggende spørsmål: Hvis sjelen er en ikke-utstrakt substans, hvordan kan den forårsake frivillige handlinger ved å påvirke kroppens ånder? Descartes svarer 21. mai at spørsmålet må forstås på en annen måte enn gjennom mekaniske forklaringer. 

Problemstillingen i denne teksten er derfor: \textit{Hvordan forsvarer Descartes at sjelen, som kun er en tenkende substans, kan styre kroppens ånder og forårsake frivillige handlinger – og hvordan viser Elisabeth svakheten i dette forsvaret?}

\section{Elisabeths innvending (6. mai 1643)}

Elisabeth skriver at hun ikke kan forstå hvordan sjelen, som ifølge Descartes er uten utstrekning, kan virke på kroppen. Hun viser til at bevegelse krever berøring eller fysisk påvirkning. Som hun formulerer det, berøring synes meg uforenlig med en immateriell substans (Cappelen, Torsen, Watzl, 2021, s. 216). 

For Elisabeth er dette et logisk problem: I følge datidens mekaniske fysikk må alt som beveger noe annet, gjøre det ved kontakt, kraft eller trykk. Hvis sjelen ikke har utstrekning, kan den ikke berøre kroppen, og dermed heller ikke forårsake bevegelse. 

Hun ber derfor Descartes gi en mer presis definisjon av sjelen – ikke bare som en substans som \textit{tenker}, men som noe som også kan forklares i forhold til sine virkninger. Elisabeth anerkjenner at det kan finnes sjeler uten tenkning, som hos barn eller besvimte personer, og ønsker derfor en forklaring på hva sjelen \textit{er} i seg selv, ikke bare hva den gjør.

\section{Descartes’ svar (21. mai 1643)}

Descartes innleder sitt svar med å rose Elisabeths spørsmål som et av de mest presise som kan stilles. Han erkjenner at han hittil har fokusert på å forklare sjelen som tenkende substans, men ikke dens forening med kroppen. Han beskriver at kunnskapen om menneskesjelen hviler på to forhold: at den tenker, og at den er forent med kroppen slik at den kan virke og påvirkes sammen med den (s. 217).

Deretter understreker Descartes at all sikker viten krever at vi holder ideene om sjel og kropp atskilt. Det er, sier han, en feil å forsøke å forklare et fenomen som tilhører sjelen ved hjelp av begreper som gjelder kroppen, eller omvendt. Elisabeth forsøker å forstå sjelens virkning på kroppen ut fra mekaniske prinsipper som gjelder for materielle ting, men dette er ifølge Descartes en feil anvendelse av begreper. 

Når vi forsøker å forklare forholdet mellom sjel og kropp mekanisk, gjør vi samme feil som når vi prøver å forstå tyngdekraften som et fysisk trykk: vi vet at den virker, men ikke hvordan. På samme måte opplever vi at viljen beveger kroppen, uten at vi kan forklare den årsaksmessige forbindelsen. Dette er en form for \textit{grunnleggende erfaring}, ikke et fysisk prinsipp (s. 218).

Slik flytter Descartes problemet fra en mekanisk til en erfaringsmessig forklaring. Han hevder at vi har en medfødt idé om forbindelsen mellom vilje og handling, og at dette er tilstrekkelig til å forstå hvordan sjelen kan bevege kroppen. Å kreve en mekanisk forklaring er derfor å stille et galt spørsmål.

\section{Drøfting}

Elisabeths spørsmål tvinger Descartes til å gå utover sin tidligere forståelse av sjelen som bare en tenkende substans. Hennes innvending avdekker et hull i dualismen: hvis sjel og kropp er helt ulike, hvordan kan de påvirke hverandre? Hun bruker datidens fysikk som et logisk rammeverk, og hennes spørsmål uttrykker et krav om sammenheng og årsaksforklaring.

Descartes sitt svar er å avvise at forholdet mellom sjel og kropp kan forstås mekanisk. For ham er foreningen mellom sjel og kropp en grunnleggende erfaring, på samme måte som vi intuitivt forstår tyngdens virkning. Dette er en viktig nyanse i hans filosofi: det viser at mennesket må forstås både som tenkende og handlende, uten at det ene kan reduseres til det andre.

Men Descartes unngår samtidig å gi en faktisk forklaring på hvordan sjelen påvirker kroppen. Han erkjenner at vi ikke kan beskrive dette gjennom årsak og virkning slik vi gjør i fysikken, men han gir heller ingen alternativ mekanisme. Dermed står Elisabeths innvending fremdeles som gyldig: hun etterspør en forklaring på hvordan handling oppstår, ikke bare en påpekning av at den gjør det.

Descartes’ forsvar innebærer derfor et skifte i perspektiv – fra forklaring til erfaring. Foreningen mellom sjel og kropp må erfares direkte gjennom vår egen vilje. Likevel kan man hevde at dette er utilfredsstillende som filosofisk svar, siden det hviler på en forutsetning som ikke lar seg etterprøve.

\section{Konklusjon}

I brevet 21. mai 1643 forsvarer Descartes sin påstand om at sjelen kan styre kroppen ved å vise til menneskets erfaring av å ville og handle. Han hevder at forbindelsen mellom sjel og kropp ikke kan forstås mekanisk, men bare gjennom den umiddelbare erfaringen av vår egen natur. Elisabeths brev 6. mai avslører likevel svakheten i dette synet: Descartes gir ingen forklaring på hvordan en ikke-utstrakt sjel faktisk kan påvirke en utstrakt kropp. Hennes spørsmål blir derfor stående som et av de mest presise og varige problemene for dualismen.

\section{Referanser}

\begin{itemize}
    \item Cappelen, H., Torsen, I. \& Watzl, S. (2021). \textit{Vite, være, gjøre: Exphil}. Gyldendal. (s. 216–219)
\end{itemize}

\end{document}
